\chapter{Onde achar conteúdo}
\label{ap:onde-achar}

Tudo o que a apostila cita, separado por tipo, com o lugar onde comprar ou
baixar cada obra. As colunas saem dos campos do \texttt{.bib} (convenção no
cabeçalho de \texttt{latex/bib/testes.bib}): nenhum endereço foi copiado à mão.

\begin{description}
\item[Comprar] Página de venda da própria editora. Preferimos a editora à
  loja; loja só quando a editora não vende direto. \enquote{Esgotado} quando
  a editora não vende mais.
\item[Edição PT] Edição brasileira, com link para a editora brasileira quando
  ela vende o livro. Obras escritas em português aparecem como
  \enquote{original em português} e levam a marca \pt.
\item[Download gratuito] Cópia legal e gratuita: livro aberto, PDF na página
  do autor ou da conferência, repositório da universidade, arXiv ou o DOI
  quando a editora abriu o acesso. Nunca cópias não oficiais.
\end{description}

\nocite{livro:valente2020}
\nocite{curso:dev-java-learn}
\nocite{comunidade:soujava,comunidade:soujava-github,comunidade:guj}

\section{Livros}
Para começar em português: \citetitle{livro:aniche2015} \pt e
\citetitle{livro:valente2020} \pt, este com leitura gratuita no site do
autor.

\ondeachar{livro}

\section{Papers e teses}
Papers não se compram avulsos aqui: o link leva à cópia aberta (página do
autor ou da conferência, repositório da universidade, arXiv) ou ao DOI. Os
papers da ACM saem pelo DOI: desde 1º de janeiro de 2026 toda a ACM Digital
Library é de acesso aberto. Sem cópia aberta, o link é o DOI, marcado
\enquote{acesso pago}.

\ondeachar*{paper}

\section{Vídeos}
A maioria em inglês; os marcados com \pt são em português. Todos gratuitos.

\ondeachar*{video}

\section{Artigos e blog posts}

\ondeachar*{artigo}

\section{Cursos gratuitos}
Nenhum curso pago é indicado. Além dos tutoriais abaixo, os guias oficiais do
Quarkus e do GraalVM citados nas Partes~III e~IV funcionam como curso
gratuito (lista completa na bibliografia, seção de documentação).

\ondeachar*{curso}

\section{Comunidades em português}
Lugares para perguntar, achar eventos e conhecer quem trabalha com Java no
Brasil.

\ondeachar*{comunidade}

\section{Documentação e ferramentas}
Os sites oficiais das ferramentas estão no apêndice~\ref{ap:ferramentas};
especificações, JEPs e guias, na bibliografia (apêndice~\ref{ap:bibliografia}).
